%==============================================================================
\chapter{Dictionnaire de données et règles de calcul}
%==============================================================================
Ce chapitre répond à la question des inputs et des outputs du système. Il définit
les données manipulées, leurs types et les formules de calcul des résultats.

\section{Conventions de types}
Les types utilisés sont : entier, décimal, texte, date, heure, datetime, booléen,
énumération et référence (clé étrangère vers une autre entité). Les grandeurs
physiques sont associées à une unité paramétrable pour permettre
l'internationalisation des systèmes de mesure.

\section{Dictionnaire des données (inputs)}
\begin{xltabular}{\textwidth}{>{\bfseries}p{2.6cm} p{3.4cm} p{2.4cm} X}
\toprule
\textbf{Entité} & \textbf{Attribut} & \textbf{Type} & \textbf{Description} \\
\midrule
\endhead
Fermier & id, nom, contact & entier, texte & Propriétaire exploitant. \\
Ferme & id, nom, fermier & entier, texte, référence & Exploitation rattachée à un fermier. \\
Site & id, nom & entier, texte & Emplacement d'apiculture. \\
Site & latitude, longitude & décimal (degrés) & Géolocalisation. \\
Site & altitude & décimal (unité paramétrable) & Altitude du site. \\
Site & dateMiseEnOeuvre, dateDemenagement, dateCloture & date & Cycle de vie du site, deux dernières optionnelles. \\
Ruche & id, modele & entier, texte & Ruche hébergée par un site. \\
Ruche & nbHausses & entier (0 à 5) & Nombre d'étages d'extension. \\
Ruche & ferme & référence & Ferme propriétaire (distincte du site d'emplacement). \\
Ruche & etat & énumération & État du cycle de vie (créée, peuplée, active, en division, en collecte, clôturée), voir figure~\ref{fig:etat}. \\
Ruche & agentResponsable & référence & Agent responsable à l'instant donné. \\
Compartiment & id, type & entier, énumération & Compartiment de la ruche : corps ou hausse. \\
Compartiment & capaciteCadres & entier (1 à 10) & Nombre de slots de cadres. \\
Cadre & id, slot, type & entier, énumération & Cadre installé sur un slot identifié. \\
Cadre & poids & décimal (unité paramétrable) & Poids du cadre, base du calcul de la quantité de miel. \\
Agent & id, nom, role & entier, texte, énumération & Apiculteur, superviseur, responsable ou administrateur. \\
Planning & id, ruche, agent & entier, référence & Planning de visite. \\
Planning & datePrevue, statutApprobation & date, énumération & Date prévue et approbation du superviseur. \\
Planning & superviseur & référence & Agent superviseur qui approuve le planning. \\
Visite & id, ruche, agent & entier, référence & Visite réalisée. \\
Visite & planning & référence & Planning approuvé à l'origine de la visite (optionnel). \\
Visite & date, heure, duree & date, heure, entier (min) & Horodatage et durée. \\
Visite & raison & énumération & Motif de la visite. \\
Visite & constatations, actionsPrevues, actionsEffectuees, recommandations & texte & Contenu du rapport. \\
Visite & effectifQualitatif, etatSante & énumération & Évaluation de l'essaim. \\
Visite & productivite & entier (1 à 3) & Niveau de productivité. \\
Photo & id, url, visite & entier, texte, référence & Photo attachée à un rapport de visite. \\
Mesure & id, ruche, typeIndicateur & entier, référence, énumération & Mesure d'un indicateur. \\
Mesure & valeur, unite & décimal, texte & Valeur et unité paramétrable. \\
Mesure & horodatage, latitude, longitude & datetime, décimal & Mesure horodatée et géolocalisée. \\
Récolte & id, ruche, date & entier, référence, date & Récolte de miel. \\
Récolte & quantiteMiel, lotQR & décimal (unité paramétrable), texte & Quantité et identifiant de lot. \\
Seuil & id, type, valeur, unite & entier, énumération, décimal, texte & Seuil paramétrable du système. \\
\bottomrule
\end{xltabular}

\begin{encadre}
\textbf{Schéma de référence :} le dictionnaire ci-dessus est la vue métier des
données. Sa traduction relationnelle normalisée (clés, types SQL et contraintes
\texttt{CHECK}) constitue le \textbf{modèle logique de données} (MLD), qui fait
foi pour l'implémentation et est présenté à l'annexe~\ref{chap:complements}
(figure~\ref{fig:mld}). La correspondance des noms entre le dictionnaire, le
diagramme de classes et le MLD est donnée par la table~\ref{tab:vocabulaire}.
\end{encadre}

\section{Résultats et règles de calcul (outputs)}
\begin{itemize}
  \item \textbf{Quantité de miel d'une ruche} : somme des poids des cadres de miel des hausses, diminuée de la tare des éléments.
\end{itemize}
\[ \text{QuantiteMiel}(r) = \sum_{c \in \text{cadresMiel}(r)} \text{poids}(c) - \text{tare}(r) \]
Cette valeur est exposée par le service web via \texttt{getZummHoneyActualQuantity(zummID)}.
\begin{itemize}
  \item \textbf{Retour sur investissement} d'un site ou d'une ruche :
\end{itemize}
\[ \text{ROI} = \frac{\text{Production valorisée}}{\text{Coûts des interventions}} \]
\begin{itemize}
  \item \textbf{Alerte de collecte} : déclenchée si la production dépasse le seuil paramétrable.
\end{itemize}
\[ \text{Production}(r) > \text{SeuilProduction} \Rightarrow \text{alerte collecte ou extension} \]
\begin{itemize}
  \item \textbf{Alerte de baisse} : déclenchée si la baisse relative dépasse le seuil.
\end{itemize}
\[ \frac{V_{\text{préc}} - V_{\text{actuel}}}{V_{\text{préc}}} > \text{SeuilBaisse} \Rightarrow \text{alerte inspection} \]
\begin{itemize}
  \item \textbf{Conversion d'unités} vers l'unité de référence, pour comparer des mesures hétérogènes :
\end{itemize}
\[ V_{\text{référence}} = V_{\text{mesure}} \times \text{facteurConversion}(\text{unite}) \]

\section{Requêtes types}
Les résultats s'appuient sur des requêtes exécutées via JDBC, par exemple pour la
quantité de miel d'une ruche : \texttt{SELECT SUM(quantiteMiel) FROM recolte WHERE ruche = zummID}.

